% Beamer-Klasse mit 16:9-Seitenverhältnis (32:18)
\documentclass[aspectratio=3218]{beamer}

% Pakete laden
\input{pakete}

% ICPAC-Stil laden
\usepackage{beamer_ovgu_169}

% Sprecher-Notizen (gesprochenes Skript) anzeigen? Standardmäßig aus.
% Zum Anzeigen die folgende Zeile einkommentieren:
%\setbeameroption{show notes}

\title[CRMA]{From Impact-Based Forecasting to Epistemic Risk Assessment}
\subtitle{Continuous Risk Monitoring \& Assessment (CRMA): \\ an epistemic layer for early warning and anticipatory action in East Africa}
\author{Nishadh Kalladath}
\institute[CRMA \textbar\ CRAF'd \textbar\ ICPAC]{%
	Disaster Risk Management Unit, IGAD Climate Prediction \& Applications Centre \\
	%Developed under the CRAF'd project%
}

\begin{document}

% Fußzeile (mit Logos) nur auf der Titelfolie unterdrücken
{
\setbeamertemplate{footline}{}
\begin{frame}
	\maketitle
\end{frame}
}



\begin{frame}[label=outline]{Outline}
	\vspace*{1ex}
	\begin{center}\small
	\begin{tabular}{@{}clp{9.5cm}@{}}
		\toprule
		 & \textbf{Part} & \textbf{Question it answers} \\
		\midrule
		1 & Impact-Based Forecasting & \emph{What is the risk?} - IBF as an \textbf{ontological} model \\
		2 & The missing question & \emph{How well do we know it?} - the \textbf{epistemic} gap in EWS \& anticipatory action \\
		3 & Two communities & Forecasters and DRM practitioners - and the seam between them \\
		4 & CRMA & Bayesian network + evidence curation (hard / soft / virtual) \\
		5 & Two directions & (a) AI automation of CRMA \quad (b) the DOC ``flight simulator'' \\
		\bottomrule
	\end{tabular}
	\end{center}
	\vspace*{1ex}
	\begin{block}{}
		\centering
		CRMA is an \textbf{epistemic} framework that complements the
		\textbf{ontological} assessment performed by Impact-Based Forecasting.
	\end{block}
	\note{SAY: ``The argument runs in five steps. First, impact-based forecasting
		- what it is, and what kind of question it answers. Second, the question
		it does not answer. Third, why that gap sits exactly on the seam between
		the forecasting community and the disaster-risk-management community.
		Fourth, CRMA as the layer that closes it. And fifth, the two directions
		this is now taking.''}
\end{frame}

% =====================================================================
\section{From IBF to epistemic risk}
% nur für die Miniframe-Navigation
\subsection*{}

\begin{frame}{Impact-Based Forecasting - the ontology of risk}
	IBF moved early warning from \emph{``what the weather \textbf{will be}''} to
	\emph{``what the weather \textbf{will do}''}.
	\begin{block}{The conventional risk model}
		\centering\large
		$\mathrm{Risk} \;=\; f(\;\mathrm{Hazard},\; \mathrm{Exposure},\;
		\mathrm{Vulnerability}\;)$ \\[0.6ex]
		\normalsize Impact $=$ Hazard intensity $\times$ Exposed elements
		$\times$ their Vulnerability
	\end{block}
	\begin{itemize}
		\item \textbf{Hazard} - forecast intensity, footprint, duration
			(rainfall, discharge, SPI).
		\item \textbf{Exposure} - people, crops, assets, infrastructure in the
			footprint.
		\item \textbf{Vulnerability} - the damage/loss function that converts
			exposure into impact.
	\end{itemize}
	\centering
	This is an \textbf{ontology}: a formal account of \emph{what exists in the
	world} and how those entities combine to produce loss.
	\note{SAY: ``Start where the field already is. Impact-based forecasting is
		the shift from forecasting the weather to forecasting what the weather
		will do. Its engine is the conventional risk model: risk as a function of
		hazard, exposure and vulnerability. Hazard from the forecast, exposure
		from the population and asset layers, vulnerability from a damage
		function. Notice what kind of thing this is. It is an ontology - a
		formal statement of what exists in the world and how those entities
		combine to produce loss. Every term on the right-hand side names a real
		object out there.''}
\end{frame}

\begin{frame}{What the equation answers - and what it assumes}
	\begin{columns}[T]
		\begin{column}{0.48\textwidth}
			\textbf{IBF answers well}
			\begin{itemize}
				\item \emph{What is the flood risk?}
				\item \emph{How many people may be affected?}
				\item \emph{Which admin-1 is most exposed?}
				\item \emph{What losses are expected?}
			\end{itemize}
			\vspace{1ex}
			Statements \textbf{about the world}.
		\end{column}
		\begin{column}{0.48\textwidth}
			\textbf{But it quietly assumes}
			\begin{itemize}
				\item the hazard input is \textbf{a value}, not a spread;
				\item exposure \& vulnerability layers are \textbf{current and
					correct};
				\item the damage function \textbf{transfers} to this place, this
					season;
				\item evidence arrives \textbf{once}, not continuously and in
					conflict.
			\end{itemize}
		\end{column}
	\end{columns}
	\vspace{1ex}
	\begin{block}{}
		\centering Uncertainty is \textbf{acknowledged}, but it is \textbf{not the
		object} of the assessment.
	\end{block}
	\note{SAY: ``The equation answers a specific family of questions extremely
		well - what is the risk, how many people, which region, what losses.
		These are statements about the world. But look at what it assumes to be
		able to answer them: that the hazard enters as a value rather than a
		spread; that the exposure and vulnerability layers are current and
		correct; that a damage function estimated somewhere else transfers here;
		and that evidence arrives once, cleanly, rather than continuously and in
		contradiction. Uncertainty is acknowledged in IBF - error bars,
		probabilistic thresholds - but it is never the object of the assessment.
		The object is always the risk itself.''}
\end{frame}

\begin{frame}{The missing question - epistemic risk assessment}
	\begin{columns}[T]
		\begin{column}{0.48\textwidth}
			\begin{block}{Ontological (IBF)}
				\centering \emph{What is the risk?} \\[0.8ex]
				\footnotesize The object of assessment is \\ \textbf{the flood
				risk itself}.
			\end{block}
		\end{column}
		\begin{column}{0.48\textwidth}
			\begin{block}{Epistemic (CRMA)}
				\centering \emph{Given all available evidence, what should we
				\textbf{currently believe} about the risk?} \\[0.8ex]
				\footnotesize The object of assessment is \\ \textbf{our knowledge
				of the flood risk}.
			\end{block}
		\end{column}
	\end{columns}
	\vspace{1ex}
	An epistemic assessment must therefore carry, explicitly:
	\begin{center}\small
		\textbf{confidence} \quad$\cdot$\quad \textbf{uncertainty}
		\quad$\cdot$\quad \textbf{evidence quality} \quad$\cdot$\quad
		\textbf{conflicting observations} \quad$\cdot$\quad \textbf{belief
		updating} \quad$\cdot$\quad \textbf{knowledge gaps}
	\end{center}
	\centering
	Estimating risk is not enough. We must also assess: \textbf{how justified is
	that estimate?}
	\note{SAY: ``So here is the distinction the whole talk rests on. The
		ontological question is: what is the risk? The epistemic question is:
		given all the evidence available right now, what should we believe about
		the risk? In the first, the object of assessment is the flood. In the
		second, the object of assessment is our knowledge of the flood. That
		second assessment has to carry things the first one does not: confidence,
		uncertainty, the quality of each piece of evidence, observations that
		contradict each other, how belief updates, and what we still do not know.
		Estimating the risk is not enough. We must also assess how justified that
		estimate is.''}
\end{frame}

\begin{frame}{Why this matters for EWS and anticipatory action}
	Both are built on a \textbf{trigger}: forecast $\rightarrow$ threshold
	$\rightarrow$ pre-agreed action.
	\begin{block}{}
		\centering A trigger is an \textbf{ontological device}. It asks
		\emph{``has the risk crossed X?''} - never \emph{``how well do we know
		it has?''}
	\end{block}
	\begin{itemize}
		\item \textbf{No record of belief.} We log that the trigger fired, not
			\emph{what we knew, when, and on what evidence}.
		\item \textbf{Conflicting evidence is invisible.} A gauge, a forecast and
			a field report disagree - the threshold sees only one number.
		\item \textbf{Confidence is discarded.} A 40\% exceedance and a 95\%
			exceedance both simply ``cross''.
		\item \textbf{No-regret space collapses.} Binary act/don't-act hides the
			graded response an operations centre actually uses.
		\item \textbf{Accountability is misplaced.} Decisions get judged by the
			\textbf{outcome}, not by the \textbf{evidence available at the time}.
	\end{itemize}
	\centering\footnotesize A good decision $\neq$ a good outcome - but without an
	epistemic record, we cannot tell them apart.
	\note{SAY: ``Why does this matter operationally? Because early warning
		systems and anticipatory action are both built on a trigger: forecast,
		threshold, pre-agreed action. And a trigger is an ontological device. It
		asks whether the risk has crossed a line. It never asks how well we know
		that it has. The consequences are concrete. We keep no record of belief -
		only that the trigger fired. Conflicting evidence is invisible to a
		threshold. Confidence is discarded: a forty-percent exceedance and a
		ninety-five-percent exceedance both simply cross. The no-regret space
		collapses into act or don't act. And accountability lands in the wrong
		place - decisions get judged by the outcome rather than by the evidence
		available at the time. A good decision is not the same as a good outcome,
		but without an epistemic record you cannot tell the two apart.''}
\end{frame}

\begin{frame}{Two communities, one seam}
	The gap is not technical. It is \textbf{institutional} - it falls between two
	well-established communities.
	\begin{center}\small
	\begin{tabular}{@{}lp{4.6cm}p{4.6cm}@{}}
		\toprule
		 & \textbf{Weather \& climate forecasting} & \textbf{Disaster risk management} \\
		\midrule
		\textbf{Core object} & the atmospheric state & the exposed \& vulnerable population \\
		\textbf{Native currency} & \textbf{probability} (ensembles, spread, skill) & \textbf{consequence} (loss, damage, need) \\
		\textbf{Trained to} & quantify uncertainty & act decisively despite it \\
		\textbf{Time frame} & forecast cycle (6\,h -- season) & the operational decision window \\
		\textbf{Hands over} & a probabilistic product & a plan requiring a commitment \\
		\bottomrule
	\end{tabular}
	\end{center}
	\begin{block}{The hand-off problem}
		\centering
		The forecaster's \textbf{spread} is flattened into the DRM
		actor's \textbf{decision} - and the reasoning in between is never written
		down.
	\end{block}
	\centering\footnotesize
	Each community owns half the epistemics; neither owns the \emph{record} of
	how evidence became a decision.
	\note{SAY: ``And I want to be clear that this gap is not a technical failure.
		It is institutional. It falls on the seam between two mature communities.
		The forecasting community's object is the atmospheric state; its native
		currency is probability; it is trained to quantify uncertainty honestly.
		The disaster-risk-management community's object is the exposed
		population; its currency is consequence; it is trained to act decisively
		in spite of uncertainty. Both are right within their own frame. But at
		the hand-off, the forecaster's carefully-constructed spread gets flattened
		into a single decision - and the reasoning in between is never written
		down anywhere. Each community owns half of the epistemics. Neither owns
		the record of how evidence became a decision. That record is what CRMA
		is.''}
\end{frame}

% =====================================================================
\section{CRMA}
\subsection*{}

\begin{frame}{CRMA - the epistemic layer over IBF}
	\begin{block}{}
		\centering
		Continuous Risk Monitoring \& Assessment (\textbf{CRMA}) is an
		\textbf{epistemic} framework that complements the \textbf{ontological}
		assessment performed by Impact-Based Forecasting.
	\end{block}
	CRMA is \emph{not} $\;$forecast $\rightarrow$ trigger. It is:
	\begin{block}{}
		\centering\Large
		Observe $\rightarrow$ Assess $\rightarrow$ Update $\rightarrow$ Decide
	\end{block}
	\begin{itemize}
		\item \textbf{IBF supplies the ontology} - hazard, exposure,
			vulnerability, impact. CRMA does not replace it.
		\item \textbf{CRMA supplies the epistemics} - evidence typing,
			confidence, belief updating, and an \textbf{audit trail}.
		\item Its output is a risk \textbf{indication} with its justification
			attached - \textbf{not} a calibrated probability of disaster, and
			\textbf{not} a decision.
	\end{itemize}
	\centering\footnotesize
	Two ingredients make it work: a \textbf{Bayesian network}, and
	\textbf{curated evidence} classified by type.
	\note{SAY: ``So: CRMA. Continuous Risk Monitoring and Assessment is an
		epistemic framework that complements the ontological assessment IBF
		performs. It does not replace the risk equation - it sits on top of it.
		IBF supplies the ontology; CRMA supplies the epistemics: evidence typing,
		confidence, belief updating, and an audit trail. And its shape is not
		forecast-then-trigger. It is observe, assess, update, decide - which is
		what an operations centre actually does. Two ingredients make it work: a
		Bayesian network, and evidence that has been curated and classified by
		type.''}
\end{frame}

% \begin{frame}{Scenario Simulation}
% 	\begin{block}%{In one sentence}
    
% 		An interactive scenario-based simulation that trains participants to
% 		evaluate evolving evidence, update risk assessments under uncertainty,
% 		and support operational decision-making - using real East African
% 		drought and flood events.
% 	\end{block}
% 	\begin{itemize}
% 		\item Continuous Risk Monitoring \& Assessment (\textbf{CRMA})
% 		\item Developed under the \textbf{CRAF'd} project at \textbf{ICPAC}
% 	\end{itemize}
% 	\note{SAY (0:00--0:40): ``Welcome. Some of you read the abstract for this
% 		session, From Forecast to Action. The aim of the next few minutes is to
% 		make that abstract concrete. This is a scenario-based simulation:
% 		participants step into real East African drought and flood events as
% 		they unfolded, evaluate the evidence available at the time, and work
% 		through the decisions an operations centre would face. I'll outline the
% 		idea, demonstrate it once on a single event, and then open the floor for
% 		the hands-on.''}
% \end{frame}

% \begin{frame}{What you read - the abstract}
% 	\begin{block}{}
% 		\small\itshape
% 		``This simulation immerses participants in a drought--flood escalation
% 		using the Continuous Risk Monitoring and Assessment platform developed
% 		under the CRAF'd project at ICPAC \dots\ built from \textbf{11
% 		event-based storylines} drawn from the \textbf{EM-DAT} disaster database
% 		across \textbf{11 East African countries}. Participants assume the roles
% 		of stakeholders to evaluate evolving forecast and risk evidence, and
% 		assess when anticipatory actions should be initiated through local
% 		governance structures and humanitarian actors. Decisions are explored
% 		through a \textbf{Bayesian Network} linking forecast signals, exposure,
% 		vulnerability, and action thresholds.''
% 	\end{block}
% 	\note{SAY (0:40--1:30): ``Here is the core of the abstract. A few anchors.
% 		It is event-based: eleven real disasters across eleven countries, drawn
% 		from the EM-DAT international disaster database. Participants take
% 		stakeholder roles and evaluate evidence as each event escalates,
% 		deciding when anticipatory action is justified - and through which
% 		governance and humanitarian channels. And the reasoning is organised by
% 		a Bayesian Network. I'll unpack each of these in plain terms over the
% 		next few slides.''}
% \end{frame}

\begin{frame}{Project background}
	\begin{block}{}
		\centering\small\textbf{Hazard Modeling, Impact Estimation \& Event-Based
		Climate Storylines on Drought and Flood Disasters in Eastern Africa} \\[0.5ex]
		\footnotesize\url{https://icpac-igad.github.io/e4drr/ }
        \footnotesize{Code \& docs: \url{https://github.com/icpac-igad/crma}}
	\end{block}
	\textbf{Aim:} Make the East Africa Hazard Watch (\textbf{EAHW}) a
	decision-support tool for Disaster Operation Centres (\textbf{DOC}) in their
	routine risk-monitoring \& assessment work.
	\begin{itemize}
		\item[] \textbf{Objectives}
		\item Better risk knowledge through \textbf{event-based storylines}.
		\item Impact Based Forecasting(IBF) that handles uncertainty: \textbf{Continuous Risk Monitoring
			\& Assessment}.
		\item \textbf{Capacity development} on the tools and methods from the
			first two objectives.
	\end{itemize}
	\begin{block}{}
		\textbf{Storyline} = a plausible ``what-if'' account of a hazard event. \\
		\emph{``What if the 2015/16 Super El Ni\~no returned, under today's
		exposure?''}
	\end{block}
	% \begin{center}\footnotesize
	% 	Code \& docs: \url{https://github.com/icpac-igad/crma}
	% \end{center}
	\note{SAY: ``Before the session itself, a word on the project it grows out
		of. The aim is to make the East Africa Hazard Watch a genuine
		decision-support tool for Disaster Operation Centres in their routine
		risk monitoring and assessment. Three objectives: better risk knowledge
		through event-based storylines; impact-based forecasting that handles
		uncertainty - continuous risk monitoring and assessment; and capacity
		development on the tools and methods from the first two. A storyline is
		simply a plausible what-if account of a hazard event - for example,
		what if the 2015/16 Super El Ni\~no returned under today's exposure.''}
\end{frame}

\begin{frame}{Three building blocks}
	\begin{itemize}
		\item \textbf{Decision Layer:} updates risk as new evidence arrives (Analysis of Decision/Evidence).
		\item \textbf{Data Layer:} analysis-ready, on-demand, cost-effective (API -> ARCO). \\
			\emph{(a ``Netflix / YouTube'' for risk monitoring \& assessment)}.
		\item \textbf{Risk Knowledge Layer:} storylines, Hazard Impact modeling(CLIMADA)(Contextualizing the impact).
			{\footnotesize\url{https://icpac-igad.github.io/e4drr/blog/2026-01-28-ibf-storylines/}}
	\end{itemize}
	\begin{center}
		\includegraphics[width=0.82\textwidth]{build-blocks.jpg}
	\end{center}
	\note{SAY: ``The project rests on three building blocks. The Decision Layer
		updates the risk picture as new evidence arrives. The Data Layer keeps
		everything analysis-ready, on-demand and cost-effective - think of it as
		a Netflix or YouTube for risk monitoring and assessment, streaming only
		the slice you need. And the Risk Knowledge Layer brings together the
		event storylines and local expertise.''}
\end{frame}

\begin{frame}{Decision Layer: record-keeping of an operational EWS}
	\begin{block}{}
		\centering Bayesian Network + Continuous Risk Monitoring \& Assessment
		(\textbf{CRMA})
	\end{block}
	\begin{enumerate}
		\item The best IBF system we know? \textbf{The human brain.}
		\item Forecasters already \textbf{reason conditionally under uncertainty}.
		\item \textbf{Flood or Drought fast/slow}  = a daily(every 6hr)/seasonal forecast: an initial belief,
			updated each cycle / new evidence.
	\end{enumerate}
	\textbf{Why must it be continuous?}
	\begin{enumerate}
		\item Each new ensemble cycle / observation / field report
			\textbf{changes what we know}.
		\item Record-keeping $\rightarrow$ \textbf{audit trail}: what we knew,
			when, and on what evidence.
	\end{enumerate}
	\note{SAY: ``The Decision Layer is really record-keeping for an operational
		early-warning system - a Bayesian Network paired with Continuous Risk
		Monitoring and Assessment. The framing is simple. The best
		impact-based-forecasting system we know is the human brain; forecasters
		already reason conditionally under uncertainty; and GHACOF itself is a
		seasonal forecast - an initial belief that is updated each cycle. Why
		must it be continuous? Because each new ensemble cycle, observation or
		field report changes what we know - and keeping a record of that gives
		us an audit trail: what we knew, when, and on what evidence.''}
\end{frame}

\begin{frame}{What is a Bayesian Network?}
	A \textbf{graph of evidence $\rightarrow$ risk}, with the logic written down.
	\begin{itemize}
		\item \textbf{Nodes \& links:} evidence nodes (rain, forecast, trend)
			$\rightarrow$ hidden \textbf{risk grade} $\rightarrow$ \textbf{CRMA
			decision}.
		\item \textbf{Probabilistic logic:} each link is a \textbf{conditional
			probability} - encodes ``\emph{if wet ground \& forecast deficit,
			then\dots}'' as odds, not a hard rule.
		\item \textbf{Bayesian updating:} new evidence revises the
			\textbf{posterior} - the assessment moves coherently, never ad hoc.
		\item \textbf{Record-keeper:} every state is logged - \emph{what we
			believed, when, on what evidence} - an inspectable audit trail.
	\end{itemize}
	\begin{block}{In CRMA}
		The BN \textbf{combines, updates, and makes reasoning transparent} - same
		evidence in $\rightarrow$ same assessment out. Output = a risk
		\emph{indication}, \textbf{not} a calibrated probability of disaster; the
		decision stays with people.
	\end{block}
	\note{SAY (4:00--5:00): ``The reasoning is organised by a Bayesian Network
		- and this is one of the distinctive contributions of the exercise,
		so it stays visible. But think of it as a decision-support reasoning
		framework, not a statistics lesson. It does three things: it combines
		the available evidence, it updates the assessment as conditions change,
		and it makes the reasoning transparent so you can see why the risk moved.
		What it does NOT do is make the decision. And its output is a risk
		indication - consistent and inspectable - not a validated
		probability of disaster.''}
\end{frame}



\begin{frame}{Uncertainty is the point - not the problem}
	Decisions are made \textbf{before the outcome is known} - \textbf{reason with
	uncertainty}, don't wait for certainty.
	\begin{columns}[T]
		\begin{column}{0.5\textwidth}
			\centering
			\includegraphics[height=0.46\textheight]{uncertainity.jpg} \\
			{\footnotesize \textbf{Deterministic} - one number, uncertainty
			hidden.}
		\end{column}
		\begin{column}{0.5\textwidth}
			\centering
			\includegraphics[height=0.46\textheight]{gefs.jpg} \\
			{\footnotesize \textbf{Ensemble} - a \textbf{spread}; ``40\% cross
			the threshold''. The spread \emph{is} the information.}
		\end{column}
	\end{columns}
	\centering\footnotesize The ensemble - the clearest expression of uncertainty
	- \textbf{inspires how evidence is typed} (next slide).
	\note{SAY: ``One concept underpins everything: uncertainty. In an operations
		centre the decision must be made before the outcome is known -
		uncertainty is permanent, and the skill is to reason WITH it rather than
		wait for a certainty that never arrives. The clearest place to see this
		is the forecast. A deterministic forecast gives one number and hides the
		uncertainty; an ensemble runs the model many times and exposes it as a
		spread of outcomes - and that spread is exactly the information we
		need: forty percent of members cross the flood threshold is far more
		useful than a single figure. That idea - that evidence carries
		different amounts of uncertainty - is the inspiration for how we
		classify evidence, which is the next slide.''}
\end{frame}

\begin{frame}{Evidence curation - three kinds of evidence}
	Because evidence carries \textbf{different amounts of uncertainty}, it must
	be \textbf{classified by type} - the type decides \textbf{how it enters
	the reasoning}.
	\begin{center}\small
	\begin{tabular}{@{}lp{4.0cm}p{2.9cm}p{4.4cm}@{}}
		\toprule
		\textbf{Type} & \textbf{Examples} & \textbf{In a word} & \textbf{Enters the assessment as} \\
		\midrule
		\textbf{Hard} & gauges, satellites, river levels, local knowledge & what we \textbf{measure} & a near-certain value \\
		\textbf{Soft} & forecasts, expert judgment, community reports & what we \textbf{estimate} & a probability spread (\emph{the spread is the confidence}) \\
		\textbf{Virtual} & storylines, ``what-if'' scenarios(Downward Counterfactual) & what we \textbf{imagine - to prepare} & a reliability-weighted ``likelihood'' nudge \\
		\bottomrule
	\end{tabular}
	\end{center}
	\footnotesize Classifying each piece of evidence - and judging \textbf{how
	reliable} it is - is the single most important skill the exercise trains.
	\note{SAY: ``So evidence is not all the same - precisely because it
		carries different uncertainty. HARD evidence is what we MEASURE -
		gauges, satellites, river levels, local knowledge - and it enters the
		assessment as a near-certain value. SOFT evidence is what we ESTIMATE
		- forecasts, expert judgment, community reports - uncertain, so it
		enters as a probability spread, and that spread is the confidence.
		VIRTUAL evidence is what we IMAGINE in order to prepare - storylines
		and `what-if' scenarios - injected as a reliability-weighted nudge to
		test how the picture would change. Measure, estimate, imagine. In Act
		One, participants classify every piece of evidence this way and judge how
		reliable it is - and that single act of classification decides how much
		each piece is allowed to move the assessment.''}
\end{frame}

\begin{frame}{What makes the analysis live - ARCO + streaming}
	Interactive because the data layer is \textbf{analysis-ready \&
	cloud-optimized (ARCO)}.
	\begin{itemize}
		\item \textbf{Stored once} in chunked formats (Zarr, Parquet, Icechunk) -
			read one slice without the whole archive.
		\item \textbf{Streams only the needed slice} (this date, this region) -
			updates at the \textbf{speed of conversation}, even on conference
			Wi-Fi.
		\item \textbf{Pre-computed} risk posteriors \& CRMA states per date \&
			boundary - the platform \textbf{serves a reading}, not a live model
			run.
		\item Same ARCO layers open in the \textbf{Colab / JupyterHub} notebooks.
	\end{itemize}
	\note{SAY: ``A brief word on what makes this work in a room like this. The
		decision analysis feels instant because the data layer is
		analysis-ready and cloud-optimized - ARCO. The forecasts, observations
		and risk layers are stored once in chunked, cloud-native formats, so the
		platform can read a single slice - one date, one region - without
		touching the rest of the archive. Each step streams only that slice,
		which is why the risk picture updates at the speed of conversation. The
		risk posteriors are pre-computed per date and per boundary, so the
		platform is serving a reading, not recomputing on the spot.''}
\end{frame}

% =====================================================================
\section{Two directions}
\subsection*{}

\begin{frame}{Where CRMA goes next - two directions}
	With the epistemic layer in place, the work forks into two complementary
	directions.
	\vspace{1ex}
	\begin{columns}[T]
		\begin{column}{0.48\textwidth}
			\begin{block}{(a) AI automation of CRMA}
				\emph{Machine keeps the record.} \\[0.5ex]
				Scale the evidence curation and belief updating beyond what a
				duty analyst can sustain across 11 countries, every cycle.
			\end{block}
		\end{column}
		\begin{column}{0.48\textwidth}
			\begin{block}{(b) The DOC ``flight simulator''}
				\emph{People learn the reasoning.} \\[0.5ex]
				A scenario-based simulation on 11 real East African
				drought \& flood events - the human side of the same loop.
			\end{block}
		\end{column}
	\end{columns}
	\vspace{1ex}
	\centering
	Both serve the same end: an \textbf{inspectable} risk assessment, and people
	with \textbf{calibrated trust} in it.
	\note{SAY: ``With the epistemic layer defined, the work forks. One direction
		is automation: getting an AI layer to do the evidence curation and belief
		updating at a scale no duty analyst can sustain - eleven countries,
		every forecast cycle. The other direction is human: a scenario simulator
		that trains people in exactly this reasoning. They are two halves of the
		same thing - an inspectable assessment, and people who know when to trust
		it and when to overrule it.''}
\end{frame}

\begin{frame}{Direction (a) - AI automation of CRMA}
	The bottleneck is not modelling. It is \textbf{curation}: someone must find,
	type, weigh and log each piece of evidence, every cycle.
	\begin{center}\small
	\begin{tabular}{@{}lp{5.4cm}p{5.2cm}@{}}
		\toprule
		\textbf{CRMA step} & \textbf{What the AI layer does} & \textbf{What stays human} \\
		\midrule
		\textbf{Observe} & harvest \& normalise evidence (ARCO layers, bulletins, field reports) & deciding what sources count \\
		\textbf{Curate} & propose evidence \textbf{type} (hard / soft / virtual) \& reliability & confirming or overriding the typing \\
		\textbf{Assess} & instantiate the BN, compute posteriors per admin-1 & the network structure \& priors \\
		\textbf{Update} & re-run each cycle, flag \textbf{what changed and why} & judging whether the change is credible \\
		\textbf{Decide} & draft the justification narrative \& audit entry & \textbf{the decision} \\
		\bottomrule
	\end{tabular}
	\end{center}
	\begin{block}{}
		\centering
		Automate the \textbf{record-keeping}, never the \textbf{judgement}. The
		audit trail is what makes the automation checkable.
	\end{block}
	\note{SAY: ``Direction A. The bottleneck in CRMA is not the modelling - the
		Bayesian network is cheap to run. The bottleneck is curation: somebody
		has to find each piece of evidence, decide its type, judge its
		reliability, and log it, every cycle, for every region. That is exactly
		the kind of work an AI layer can carry. It harvests and normalises the
		evidence, proposes a type and a reliability, instantiates the network,
		recomputes the posteriors, and - crucially - flags what changed and why.
		What stays human is the judgement: which sources count, whether a typing
		is right, whether a change is credible, and the decision itself. The
		principle is to automate the record-keeping and never the judgement. And
		note the epistemic framing pays for itself here: because CRMA already
		keeps an audit trail, the automation is checkable in a way an end-to-end
		black-box impact model never is.''}
\end{frame}

\begin{frame}{Direction (b) - a ``flight simulator'' for the DOC}
	\begin{itemize}
		\item \textbf{Relive} a real drought or flood - only the evidence
			available \textbf{at the time}.
		\item \textbf{Grade each admin-1} with a risk colour - the call a DOC has
			to make.
		\item \textbf{Date cursor} steps the event window: flood
			\emph{daily} ($\sim$15 d), drought \emph{monthly} (a full season).
	\end{itemize}
	\begin{block}{End outcome - a CRMA colour grade per admin-1}
		\centering
		\textcolor{green!60!black}{$\blacksquare$}~\textbf{Monitor} \quad
		\textcolor{yellow!75!black}{$\blacksquare$}~\textbf{Evaluate} \quad
		\textcolor{orange}{$\blacksquare$}~\textbf{Assess} \quad
		\textcolor{red}{$\blacksquare$}~\textbf{Actionable Risk}
	\end{block}
	\note{SAY (1:30--2:15): ``In plain terms, it is a flight simulator for an
		operations centre. Participants relive a real drought or flood as it
		unfolded, with only the information that existed at the time. Their task
		is the call a DOC actually has to make: grade every admin-1 region with a
		risk colour - green Monitor, yellow Evaluate, orange Assess, red
		Actionable Risk - and justify it on the evidence weighed. A date cursor
		steps through the event window: daily over about fifteen days for a
		flood, lead to escalation to onset; monthly over several months for a
		drought, here the OND short-rains season. The outcome is revealed only
		afterwards, and the decision is assessed against the evidence available at
		the time, not the guess.''}
\end{frame}

\begin{frame}{What this session produces}
	Beyond the learning, a concrete output:
	\begin{block}{}
		\centering
		Recommendations for co-designing tools that bridge probabilistic
		forecasting, evidence-based risk assessment, and humanitarian response.
	\end{block}
	\begin{itemize}
		\item Where did a signal speak early - and where did it mislead?
		\item Where did the assessment diverge from what happened, and why?
		\item What would make these tools usable in a real DOC?
	\end{itemize}
	\note{SAY (5:40--6:20): ``The session is also intended to produce something
		concrete: recommendations for co-designing tools that bridge
		forecasting, risk assessment, and humanitarian response. Where did a
		signal warn early, where did it mislead, where did the assessment
		diverge from what happened - and what would make these tools genuinely
		usable in an operations centre.''}
\end{frame}

\begin{frame}{"Flight" Simulator for DOC: three acts}
	The concept resolves into \textbf{three questions} a DOC asks, in order:
	\begin{center}
	\begin{tabular}{@{}cp{4.8cm}p{5.2cm}@{}}
		\toprule
		\textbf{Act} & \textbf{Focus} & \textbf{Question} \\
		\midrule
		\textbf{I} & Understanding the Event & \emph{What is happening?} \\
		\textbf{II} & Evidence Evaluation \& Risk Assessment & \emph{What do we think is happening?} \\
		\textbf{III} & Decision \& Reflection & \emph{What should we do, and why?} \\
		\bottomrule
	\end{tabular}
	\end{center}
	\centering Each act's \textbf{output is the next act's input.}
	\note{SAY (6:20--7:00): ``Everything resolves into three acts. Act One,
		Understanding the Event - what is happening. Act Two, Evidence
		Evaluation and Risk Assessment - what do we think is happening. Act
		Three, Decision and Reflection - what should we do, and why. The
		output of each act feeds the next. Let me demonstrate all three on one
		real event before participants run it themselves.''}
\end{frame}

% \begin{frame}{The evidence and decision pathway}
% 	\begin{block}{The pathway explored during the exercise}
% 		\centering
% 		Observations $\rightarrow$ Forecast Evidence $\rightarrow$ Risk
% 		Assessment $\rightarrow$ Decision Options $\rightarrow$ Historical
% 		Outcome
% 	\end{block}
% 	\begin{itemize}
% 		\item \textbf{Observations \& reports} - current conditions;
% 			situational awareness.
% 		\item \textbf{Forecast evidence} - ensemble outlooks; uncertain,
% 			probabilistic signals.
% 		\item \textbf{Risk assessment} - evidence combined into a structured,
% 			consistent picture.
% 		\item \textbf{Decision options} - the operational choices a DOC must
% 			weigh.
% 		\item \textbf{Historical outcome} - the recorded loss \& damage,
% 			revealed only at the end.
% 	\end{itemize}
% 	Aligned with CRMA: \textbf{Observe $\rightarrow$ Assess $\rightarrow$ Update
% 	$\rightarrow$ Decide.}
% 	\note{SAY (2:15--3:10): ``This is the pathway participants move along:
% 		observations, forecast evidence, a structured risk assessment, the
% 		decision options, and finally the historical outcome. The storylines are
% 		constructed from real recorded events, but participants experience them
% 		forward - they observe, assess, update as evidence arrives, and
% 		decide - and only at the end is the recorded loss and damage revealed.
% 		That ordering is deliberate: it mirrors how an operations centre
% 		actually works, and it matches the CRMA philosophy - observe, assess,
% 		update, decide.''}
% \end{frame}





% \begin{frame}{What the exercise provides}
% 	\begin{center}
% 	\begin{tabular}{@{}lp{8.5cm}@{}}
% 		\toprule
% 		\textbf{Ingredient} & \textbf{In the session} \\
% 		\midrule
% 		\textbf{Ensemble forecasts} & Forecast evidence and uncertainty \\
% 		\textbf{Observations \& reports} & Current conditions and situational awareness \\
% 		\textbf{Historical risk knowledge} & Past impacts and vulnerability context \\
% 		\textbf{Risk-assessment engine} & Structured evidence-to-risk reasoning \\
% 		\textbf{Event storylines} & Real East African drought and flood events \\
% 		\bottomrule
% 	\end{tabular}
% 	\end{center}
% 	\begin{block}{}
% 		Hazard and impact modelling methods are used \textbf{behind the scenes}
% 		to construct the storylines and historical event context - they are
% 		\textbf{not required} to take part in the exercise.
% 	\end{block}
% 	\note{SAY (3:10--4:00): ``These are the materials participants work with:
% 		ensemble forecasts as the source of uncertain evidence; observations and
% 		reports for current conditions; historical risk knowledge for past
% 		impacts and vulnerability context; a risk-assessment engine that turns
% 		evidence into a consistent risk picture; and the event storylines
% 		themselves. Hazard and impact modelling sit behind the scenes - they
% 		are used to construct the storylines and the historical context, but
% 		participants do not need to engage with them to take part.''}
% \end{frame}



% \begin{frame}{Participant roles}
% 	Participants take the roles of \textbf{stakeholders around the operations
% 	table} - the DOC duty analyst, the humanitarian actor, the
% 	local-government focal point.
% 	\begin{itemize}
% 		\item \textbf{Evaluate} evolving forecast and risk evidence as a
% 			storyline escalates.
% 		\item \textbf{Assess} when anticipatory action should be initiated -
% 			through \textbf{local governance structures and humanitarian
% 			actors}.
% 		\item \textbf{Justify} the decision using the \textbf{evidence available
% 			at that moment}.
% 	\end{itemize}
% 	\note{SAY (5:00--5:40): ``Participants take stakeholder roles - the duty
% 		analyst, the humanitarian responder, the local-government focal point.
% 		As the storyline escalates, they evaluate the evidence and assess when
% 		anticipatory action should begin, and through which governance and
% 		humanitarian channels. And each decision must be defensible on the basis
% 		of the evidence available at that moment.''}
% \end{frame}





% =====================================================================
\section{Reflections}
% nur für die Miniframe-Navigation
\subsection*{}

\begin{frame}{Reflections}
	\begin{itemize}
		\item \textbf{Evidence typology} - hard / soft / virtual, and why
			forecasts are soft.
		%\item \textbf{Coherent updating} - the same evidence yields the same
		%	assessment; intuition varies, the framework does not.
		\item \textbf{Decisions under uncertainty} - cost-loss reasoning;
			\emph{a good decision $\neq$ a good outcome}.
		\item \textbf{Calibrated trust in tools} - the question is never ``is
			the assessment right?'' but ``is its reasoning transparent enough to
			know when to overrule it?''
	\end{itemize}
	\note{SAY (approx.\ 2 min): ``To close the structured part: the exercise
		trains four things - reading evidence by type; updating an assessment
		coherently; deciding under uncertainty, where a good decision and a good
		outcome are not the same; and, most valuably, calibrated trust in tools.
		The question for an operations centre is never simply whether the
		assessment is right, but whether its reasoning is transparent enough to
		know when to overrule it.''}
\end{frame}

\begin{frame}{Feedback we're seeking - join on Kahoot}
	\begin{columns}[c]
		\begin{column}{0.62\textwidth}
			\centering
			Help us \textbf{co-design better tools} - give your feedback live: \\[1.2ex]
			{\Large\textbf{\url{https://kahoot.it}}} \\[1.2ex]
			{\large Game \textbf{PIN} shared on screen now.} \\[1.5ex]
			{\footnotesize Or join the discussion: \\
			\url{https://github.com/icpac-igad/crma/discussions}}
		\end{column}
		\begin{column}{0.36\textwidth}
			\centering
			\includegraphics[height=0.66\textheight]{kahoot.jpg}
		\end{column}
	\end{columns}
	\note{SAY (approx.\ 1 min): ``Let's take the feedback on Kahoot - join at
		kahoot.it with the PIN on screen. The questions, straight from the
		abstract's aim: where did a signal warn early, and where did it mislead?
		Where did the assessment diverge from what happened? And what would make
		this usable in your own DOC or agency? These observations feed the
		verification record and the next iteration.''}
\end{frame}

% Fußzeile (mit Logo) auf der Schlussfolie unterdrücken
{
\setbeamertemplate{footline}{}
\begin{frame}
%{Take it further - thank you}
	% \begin{itemize}
	% 	\item \textbf{Run another event} - contrast a \textbf{strong-signal
	% 		drought} (\texttt{uganda\_karamoja\_drought\_2022}, $\approx$6-month
	% 		lead) with a \textbf{tail-risk flood}.
	% 	\item \textbf{Speed-run} - the same three acts in 5--10 minutes for
	% 		briefings.
	% 	\item \textbf{Go deeper} - the underlying data and methods are
	% 		available in the \textbf{Colab / JupyterHub} notebooks.
	% \end{itemize}
	\begin{center}\Large\textbf{Thank you.}\end{center}
	\vspace*{2ex}
	\begin{center}\includegraphics[height=1.4cm]{logos-t.png}\end{center}
	\note{SAY (approx.\ 30s): ``If time allows, run a second event - a
		slow-onset drought after a flash flood feels entirely different. The
		underlying data and methods are available in the notebooks for those who
		wish to go further. Thank you.''}
\end{frame}
}

% \begin{frame}{Appendix - Facilitator quick-reference}
% 	\small
% 	\begin{itemize}
% 		\item \textbf{Index:} \texttt{\dots run.app/scenario} - filter Flood /
% 			Drought.
% 		\item \textbf{Worked demo event:} \texttt{kenya\_nairobi\_flood\_2024}
% 			(tail-risk flood).
% 		\item \textbf{Strong-signal contrast:}
% 			\texttt{uganda\_karamoja\_drought\_2022} ($\approx$6-month lead).
% 		\item \textbf{Acts gate forward:} Act I needs all 13 answers; the
% 			engine/advisory is hidden until Act II.
% 		\item \textbf{Assessment is formative} - process over prediction. No
% 			leaderboard.
% 		\item \textbf{Scope note:} the hands-on centres the observe $\rightarrow$
% 			assess $\rightarrow$ update $\rightarrow$ decide pathway. Hazard and
% 			impact modelling are used behind the scenes (not required);
% 			\textbf{recorded EM-DAT loss is the ground-truth reference} in the
% 			debrief.
% 		\item \textbf{Offline:} questionnaire, decision \& debrief work without
% 			backend; the live engine map / DAG / advisory need \texttt{crma-api}.
% 	\end{itemize}
% 	\note{NOTE: facilitator-only - skip past in delivery.}
% \end{frame}

\end{document}
